\section{Proof of Theorem~\ref{thm:main}}\label{sec:end}
\begin{proposition}\label{lem:tail}
Suppose each $\eta$ admits a pair $\sigma,\kappa$ in
Table~\ref{tab:parameters} with the following property. Both exterior
bounds hold for ordered triples satisfying \eqref{eq:high} for sufficiently
large $x$. Suppose also that for each admissible $\eta,\sigma,\kappa$ some
$\mathcal B_X$ is empty. Then for each such $\eta$ there is a threshold
$X_\eta$ with
\[
                 H<R^{1/3-\eta}\qquad (x\ge X_\eta)
\]
for every triple.
\end{proposition}
\begin{proof}
Fix the pair in the first hypothesis and take the larger of the exterior
and exclusion thresholds for that pair. Exchanging $a,b$ preserves $x,H$.
An ordered triple above this threshold satisfying \eqref{eq:high} would
belong to the empty set $\mathcal B_X$.
\end{proof}
\begin{lemma}\label{lem:assembly}
A uniform polynomial bound and the two hypotheses of
Proposition~\ref{lem:tail} imply Theorem~\ref{thm:main}.
\end{lemma}
\begin{proof}
Given $\eps>0$, choose $\eta$ as in Section~\ref{sec:cap} with
$1/9-\eta<\eps$. For $x\ge X_\eta$ the last proposition gives
$H<R^{1/3-\eta}\le R^{2/9+\eps}$. For $x<X_\eta$ use
$H\le C_0R^k\le C_0\exp(3kX_\eta)$. Thus one may take
$C_\eps=\max(1,C_0\exp(3kX_\eta))$.
The exterior thresholds use the constants of Fact~\ref{fact:local}.
The polynomial bound supplies $(C_0,k)$ independently.
\end{proof}
\begin{proof}[Proof of Theorem~\ref{thm:main}]
Proposition~\ref{lem:cap} gives the polynomial bound. Use the pair
in Proposition~\ref{lem:small} and apply Proposition~\ref{lem:real}
at this same $\kappa$. Lemma~\ref{lem:empty} gives the exclusion.
Lemma~\ref{lem:assembly} concludes the proof.
\end{proof}
